\documentclass{subfiles}
\begin{document}\label{sec:integersEncoding}
    Many applications use integers to represent data. For instance,
        when we transmit an image, what we actualy transmit are triplets of 
        integers in the range [0, 255]. Additionaly, 
        many theoretical results we shall see in the following,
        make use of integers (see \Cref{sec:burrowsWheeler} and \Cref{sec:lempelZiv}).

    Therefore the questions: can we define a optimal code for integers?
        Can we use the codes seen so far? Both questions have an affirmative answer,
        but before analyzing it, we formalize the problem.
        
    Let \(S = s_{1}s_{2}\ldots\) be a sequence of (positive) integers.
        Is there a prefix code such that the average codeword length is 
        \(\bigO{\log s_{i}}\), for any integer \(s_{i}\).

    In the subsequent sections we describe three of such code,
        while we use the remainder of the present section to make some remarks.

    \begin{remark*}
        We can relax the constrain on positive integers. In fact, 
            we can simply proceed to define a bijective functions 
            that maps each integer to another. 
            One of such map is the following:
            \[
                f(n) = \begin{cases}
                    2n, \text{if } n > 0 \\
                    -2n + 1, \text{otherwise.} 
                \end{cases}
            \]
    \end{remark*}

    \begin{remark*}
        Why don't we use the binary representation? 
            It easily seen that this kind of representation is not prefix free.
            For instance, if we consider the binary representation of 2 and 4
            we have, respectively, 10 and 100 which are prefix one of another.
    \end{remark*}

    A simple solution is the unary encoding. 
        Essentially, any positive integer \(n\) is encoded as a sequence of \(n - 1\) zeros 
        followed by a one. That is 
        \[
            U(n) = \underbrace{00\ldots0}_{n - 1 \text{ times}}1.
        \]

    The code is trivially prefix-free, as any two different integers 
        differ by at least one digit. At the same time its use isn't feasable,
        as the number of bits needed to encode the integer increases linearly.

    Before we describe any particular family of codes to encode integers, 
        let us consider Huffman encoding. Can we use it?
        It's known by now that Huffman encoding performs well when the probabilities 
        are known a priori. Hence, if the distribution of integers is known, 
        Huffman encoding is still worth considering. On the other hand,
        if the distribution is unknown, or it may vary, then Huffman works poorly
        and special codes should be considered.

    \subsection{Elias codes}
    \subfile{../sub/4.1 - Elias codes}

    \subsection{Fibonacci codes}
    \subfile{../sub/4.2 - Fibonacci codes}
\end{document}
